\section{Edge Case Tests}

This file tests boundary conditions and edge cases from the Scriba reference.

%% ==========================================================
%% EDGE CASE A: Animation with exactly 1 step (minimum)
%% ==========================================================
\subsection{A: Single-Step Animation}

\begin{animation}[id="edge-a-single-step", label="Minimum 1-step animation"]
\shape{a1}{Array}{size=3, data=[10,20,30], label="one step"}

\step
\narrate{This animation has exactly one step --- the minimum.}
\end{animation}

%% ==========================================================
%% EDGE CASE B: Animation with 30 steps (soft limit boundary)
%% ==========================================================
\subsection{B: 30-Step Animation (Soft Limit)}

\begin{animation}[id="edge-b-30-steps", label="Exactly 30 frames"]
\shape{b}{Array}{size=5, data=[0,0,0,0,0], labels="0..4", label="$counter$"}

\step
\apply{b.cell[0]}{value=1}
\narrate{Step 1 of 30.}

\step
\apply{b.cell[1]}{value=2}
\narrate{Step 2 of 30.}

\step
\apply{b.cell[2]}{value=3}
\narrate{Step 3 of 30.}

\step
\apply{b.cell[3]}{value=4}
\narrate{Step 4 of 30.}

\step
\apply{b.cell[4]}{value=5}
\narrate{Step 5 of 30.}

\step
\recolor{b.cell[0]}{state=done}
\narrate{Step 6 of 30.}

\step
\recolor{b.cell[1]}{state=done}
\narrate{Step 7 of 30.}

\step
\recolor{b.cell[2]}{state=done}
\narrate{Step 8 of 30.}

\step
\recolor{b.cell[3]}{state=done}
\narrate{Step 9 of 30.}

\step
\recolor{b.cell[4]}{state=done}
\narrate{Step 10 of 30.}

\step
\recolor{b.cell[0]}{state=current}
\narrate{Step 11 of 30.}

\step
\recolor{b.cell[1]}{state=current}
\narrate{Step 12 of 30.}

\step
\recolor{b.cell[2]}{state=current}
\narrate{Step 13 of 30.}

\step
\recolor{b.cell[3]}{state=current}
\narrate{Step 14 of 30.}

\step
\recolor{b.cell[4]}{state=current}
\narrate{Step 15 of 30.}

\step
\recolor{b.cell[0]}{state=idle}
\narrate{Step 16 of 30.}

\step
\recolor{b.cell[1]}{state=idle}
\narrate{Step 17 of 30.}

\step
\recolor{b.cell[2]}{state=idle}
\narrate{Step 18 of 30.}

\step
\recolor{b.cell[3]}{state=idle}
\narrate{Step 19 of 30.}

\step
\recolor{b.cell[4]}{state=idle}
\narrate{Step 20 of 30.}

\step
\apply{b.cell[0]}{value=21}
\narrate{Step 21 of 30.}

\step
\apply{b.cell[1]}{value=22}
\narrate{Step 22 of 30.}

\step
\apply{b.cell[2]}{value=23}
\narrate{Step 23 of 30.}

\step
\apply{b.cell[3]}{value=24}
\narrate{Step 24 of 30.}

\step
\apply{b.cell[4]}{value=25}
\narrate{Step 25 of 30.}

\step
\recolor{b.cell[0]}{state=good}
\narrate{Step 26 of 30.}

\step
\recolor{b.cell[1]}{state=good}
\narrate{Step 27 of 30.}

\step
\recolor{b.cell[2]}{state=good}
\narrate{Step 28 of 30.}

\step
\recolor{b.cell[3]}{state=good}
\narrate{Step 29 of 30.}

\step
\recolor{b.cell[4]}{state=good}
\narrate{Step 30 of 30 --- exactly at the soft limit.}
\end{animation}

%% ==========================================================
%% EDGE CASE C: Empty narrate
%% ==========================================================
\subsection{C: Empty Narrate}

\begin{animation}[id="edge-c-empty-narrate", label="Empty narrate test"]
\shape{c}{Array}{size=2, data=[1,2], label="empty narrate"}

\step
\narrate{}

\step
\narrate{This step has narration, the previous did not.}
\end{animation}

%% ==========================================================
%% EDGE CASE D: Multiple primitives of same type
%% ==========================================================
\subsection{D: Two Arrays and Two Graphs}

\begin{animation}[id="edge-d-same-types", label="Duplicate primitive types"]
\shape{arr1}{Array}{size=3, data=[1,2,3], label="$A_1$"}
\shape{arr2}{Array}{size=4, data=[10,20,30,40], label="$A_2$"}
\shape{g1}{Graph}{nodes=["X","Y"], edges=[("X","Y")], directed=true, label="$G_1$"}
\shape{g2}{Graph}{nodes=["P","Q","R"], edges=[("P","Q"),("Q","R")], directed=false, label="$G_2$"}

\step
\recolor{arr1.cell[0]}{state=current}
\recolor{arr2.cell[3]}{state=good}
\recolor{g1.node[X]}{state=current}
\recolor{g2.node[R]}{state=done}
\narrate{Operating on two Arrays and two Graphs simultaneously.}

\step
\recolor{arr1.cell[2]}{state=done}
\recolor{arr2.cell[0]}{state=error}
\recolor{g1.edge[(X,Y)]}{state=good}
\recolor{g2.edge[(P,Q)]}{state=dim}
\narrate{Each primitive tracks its own state independently.}
\end{animation}

%% ==========================================================
%% EDGE CASE E: Very long narrate with lots of inline math
%% ==========================================================
\subsection{E: Long Narrate with Heavy Math}

\begin{animation}[id="edge-e-long-narrate", label="Long narrate text"]
\shape{e}{Array}{size=3, data=[1,2,3], label="$arr$"}

\step
\narrate{Consider the recurrence $T(n) = 2T(n/2) + O(n)$ which by the Master Theorem yields $T(n) = O(n \log n)$. For dynamic programming, the transition is $dp[i][j] = \min_{k=i}^{j-1}(dp[i][k] + dp[k+1][j] + \sum_{m=i}^{j} a_m)$. The expected value is $E[X] = \sum_{x} x \cdot P(X = x) = \int_{-\infty}^{\infty} x f(x) dx$ and the variance satisfies $\text{Var}(X) = E[X^2] - (E[X])^2 \geq 0$. By the Cauchy--Schwarz inequality, $\left(\sum_{i=1}^{n} a_i b_i\right)^2 \leq \left(\sum_{i=1}^{n} a_i^2\right)\left(\sum_{i=1}^{n} b_i^2\right)$. Furthermore, $\binom{n}{k} = \frac{n!}{k!(n-k)!}$ and $\sum_{k=0}^{n}\binom{n}{k} = 2^n$.}

\step
\recolor{e.cell[1]}{state=current}
\narrate{After all that math, we highlight cell 1 with $O(1)$ cost.}
\end{animation}

%% ==========================================================
%% EDGE CASE F: Array with size=1 (minimum)
%% ==========================================================
\subsection{F: Size-1 Array}

\begin{animation}[id="edge-f-size1-array", label="Minimum size array"]
\shape{f}{Array}{size=1, data=[42], label="singleton"}

\step
\recolor{f.cell[0]}{state=current}
\narrate{An array with exactly one element: $a[0] = 42$.}

\step
\recolor{f.cell[0]}{state=done}
\apply{f.cell[0]}{value=99}
\narrate{Updated the single cell to 99.}
\end{animation}

%% ==========================================================
%% EDGE CASE G: Graph with single node, no edges
%% ==========================================================
\subsection{G: Single-Node Graph}

\begin{animation}[id="edge-g-single-node", label="Isolated node graph"]
\shape{iso}{Graph}{nodes=["alone"], edges=[], directed=false}

\step
\recolor{iso.node[alone]}{state=current}
\narrate{A graph with a single node and zero edges.}

\step
\recolor{iso.node[alone]}{state=done}
\narrate{Mark the isolated node as done.}
\end{animation}

%% ==========================================================
%% EDGE CASE H: Highlight ephemeral clearing across steps
%% ==========================================================
\subsection{H: Highlight Ephemeral Clearing}

\begin{animation}[id="edge-h-ephemeral", label="Highlight auto-clear test"]
\shape{h}{Array}{size=4, data=[10,20,30,40], labels="0..3", label="ephemeral"}

\step
\highlight{h.cell[1]}
\narrate{Cell 1 is highlighted (ephemeral). It should auto-clear on the next step.}

\step
\narrate{No commands here --- cell 1 should no longer be highlighted.}

\step
\highlight{h.cell[3]}
\recolor{h.cell[0]}{state=done}
\narrate{Cell 3 highlighted (ephemeral), cell 0 recolored (persistent). Next step should clear only the highlight.}

\step
\narrate{Cell 3 highlight should be gone. Cell 0 should still be done (persistent).}
\end{animation}

%% ==========================================================
%% EDGE CASE I: Multiple annotate on same target, different arrow_from
%% ==========================================================
\subsection{I: Multiple Annotations Same Target}

\begin{animation}[id="edge-i-multi-annotate", label="Multiple annotations"]
\shape{dp}{DPTable}{n=5, label="dp", labels="0..4"}

\step
\apply{dp.cell[0]}{value=0}
\recolor{dp.cell[0]}{state=done}
\apply{dp.cell[1]}{value=3}
\recolor{dp.cell[1]}{state=done}
\apply{dp.cell[2]}{value=7}
\recolor{dp.cell[2]}{state=done}
\narrate{Base cases filled: $dp[0]=0$, $dp[1]=3$, $dp[2]=7$.}

\step
\recolor{dp.cell[3]}{state=current}
\annotate{dp.cell[3]}{label="+4", arrow_from="dp.cell[1]", color=good}
\annotate{dp.cell[3]}{label="+1", arrow_from="dp.cell[2]", color=info}
\narrate{Two transition arrows into $dp[3]$: one from $dp[1]$ (cost +4) and one from $dp[2]$ (cost +1).}

\step
\apply{dp.cell[3]}{value=8}
\recolor{dp.cell[3]}{state=done}
\narrate{Chose $dp[3] = \min(3+4, 7+1) = 7$... wait, $dp[3] = 8$. Annotations persist.}
\end{animation}

%% ==========================================================
%% EDGE CASE J: Deeply nested foreach (2 levels)
%% ==========================================================
\subsection{J: Nested Foreach (2 Levels)}

\begin{animation}[id="edge-j-nested-foreach", label="2-level foreach nesting"]
\shape{grid}{Grid}{rows=3, cols=3, data=[0,0,0,0,0,0,0,0,0], label="$M$"}

\step
\narrate{A $3 \times 3$ grid, all zeros. We will iterate with nested foreach.}

\step
\foreach{r}{0..2}
  \foreach{c}{0..2}
    \recolor{grid.cell[${r}][${c}]}{state=done}
  \endforeach
\endforeach
\narrate{Nested foreach: marked all 9 cells as done in one step.}

\step
\foreach{i}{0..2}
  \recolor{grid.cell[${i}][${i}]}{state=current}
\endforeach
\narrate{Diagonal cells highlighted via single-level foreach after the nested one.}
\end{animation}
